\documentclass[10pt,twocolumn]{article}
\usepackage[letterpaper,margin=0.72in,columnsep=0.24in]{geometry}
\usepackage{times}
\usepackage{microtype}
\usepackage{booktabs}
\usepackage{tabularx}
\usepackage{graphicx}
\usepackage{xcolor}
\usepackage{enumitem}
\usepackage{amsmath}
\usepackage{url}
\usepackage[hidelinks]{hyperref}
\newcommand{\system}{\textsc{ECPA}}
\newcommand{\tbd}{\textcolor{red!70!black}{\textbf{TBD}}}
\newcommand{\code}[1]{\texttt{#1}}
\setlist{nosep,leftmargin=*}
\raggedbottom
\setlength{\emergencystretch}{2em}
\def\UrlBreaks{\do\/\do-\do\_}
\sloppy
\hbadness=10000
\title{Evidence-Carrying Plugin Architecture for LLM Inference Systems}
\author{Anonymous Authors}
\date{}

\begin{document}
\maketitle

\begin{abstract}
LLM inference engines increasingly expose plugins for schedulers, KV-cache
connectors, policies, and control-plane integrations.  Existing plugin loaders
answer whether a package can be found and imported, but operators need a
stronger answer: did a compatible implementation acquire the intended
capability in every relevant process, without conflicting with another plugin,
and can the system return to a known-good configuration?  This distinction is
easy to lose under version skew, heterogeneous providers, multi-process
loading, and external services.  We present \system, an evidence-carrying
architecture that separates typed static discovery, capability negotiation,
conflict planning, saved activation intent, and process-owned runtime evidence.
Its central invariant is that no lifecycle state may claim more than its
evidence proves: installed, configured, enabled, and runtime-effective are
independent predicates.  Activation fails closed on missing or contradictory
evidence, while deterministic plans retain an explicit rollback target and
providers never inherit authority over the runtime or external services.

This paper is an implementation-grounded design and evaluation draft.  The
current prototype validates manifests and plans across several provider types,
passes 66 CPU tests, and exposes a concrete namespace regression in which all
unit tests pass while a real installed plugin becomes undiscoverable.  Formal
multi-process, composability, version-skew, rollback, and overhead campaigns
are preregistered but not yet complete; their results are marked \tbd{} rather
than inferred from existing integration evidence.  The proposed evaluation
tests whether evidence-carrying activation eliminates false-effective claims,
detects resource conflicts, survives upgrades and rollback, and remains off
the request critical path at bounded cost.
\end{abstract}

\section{Introduction}
Plugins are becoming the de facto change boundary for fast-moving inference
systems.  A scheduler policy may be delivered as a Python entry point; a KV
connector may depend on an external service; a production-stack adapter may
translate intent into Kubernetes configuration; and a provider may map a
portable capability to a host-specific command-line flag.  This modularity is
valuable, but it creates a gap between packaging facts and runtime facts.

Consider an operator who installs a plugin distribution, sees its configuration
saved, enables it, and launches an engine with multiple worker processes.  The
parent process may discover the distribution while one worker imports a
different version, never invokes the hook, or lacks the provider.  Two plugins
may independently validate yet both claim the same KV connector option.  An
external service may fail after intent is persisted.  Conventional loaders
often collapse these cases into an ``enabled'' bit.  That bit is useful intent,
but it is not evidence that the plugin changed the running system.

The repository that motivated this work contains a smaller but revealing
incident.  Its documented static bundle namespace is
\code{vllm\_hust.extension\_bundles}.  A one-line ``bugfix'' changed discovery
to \code{vllm.extension\_bundles}.  Sixty-six unit tests still passed.  Yet an
installed Mooncake example registered one bundle in the documented namespace,
the modified manager selected zero entries, and discovery returned no bundles.
Installation succeeded; the code imported; runtime intent could still exist;
the plugin was nevertheless absent from the activation path.  The incident is
not important because of the spelling error itself.  It demonstrates that
tests and static state can agree while authority and runtime effect disagree.

We ask: \emph{how can a plugin architecture make extensions discoverable,
composable, verifiable, and rollback-safe when runtime effectiveness cannot be
proved statically?}  Our answer is \system, an evidence-carrying activation
architecture.  A decision is not ``enable plugin $p$'' but an activation plan
for a concrete host version, provider set, plugin set, process topology, and
resource-ownership graph.  The plan records evidence obligations.  Processes
that actually invoke implementations satisfy those obligations; a manager
cannot manufacture runtime effectiveness from configuration.

This paper makes four intended contributions, with claim strength bounded by
the evidence available today:
\begin{itemize}
  \item a typed manifest and capability model that keeps discovery static and
  separates project namespaces from delegated upstream hooks;
  \item a compositional planner that negotiates capabilities and detects
  multiply owned resources before launch;
  \item an explicit lifecycle and evidence protocol in which process-owned
  attestations, fail-closed activation, and deterministic rollback prevent
  false-effective state; and
  \item a provider/runtime authority model and falsifiable evaluation contract
  spanning correctness, conflicts, skew, recovery, multi-process consistency,
  and overhead.
\end{itemize}
The first and parts of the second are implemented.  Process-owned evidence and
the full formal evaluation remain work in progress; we do not present the
existing provider demonstrations as proof of the full paper claim.

\section{Motivation and Problem}
\subsection{Why Existing State Is Ambiguous}
A package manager proves that files exist.  Entry-point discovery proves that
metadata contains a name.  Configuration proves that an operator supplied
parameters.  An enabled bit proves desired intent.  None proves that the
selected implementation executed in every process that matters.  Conversely,
an external service can become unreachable while enabled intent should remain
for diagnosis and later recovery.  A linear state machine both overclaims
success and destroys useful failure context.

Vanilla entry points provide a deliberately lightweight extension mechanism
\cite{vllm-plugin}.  Manual integration can add flags, environment variables,
and operator checks.  Both are legitimate baselines.  Their limitation is not
that they can never be made safe; it is that evidence obligations, ownership,
and rollback are convention rather than a common machine-checkable contract.

\subsection{System Model}
We model a host runtime $H$, providers $P$, plugins $G$, processes $W$, and
resources $R$.  A plugin manifest declares required and provided capabilities,
version predicates, claimed resources, provider needs, and evidence
obligations.  A provider translates an accepted plan into host-specific
configuration, but it does not supervise the host or acquire authority over an
external service.  Each launch has a unique identity $\ell$ and an expected
process set $W_\ell$.

An activation plan $A=(H,P,G,W,R,V)$ is admissible only if capability
requirements resolve, version predicates hold, resource ownership is
conflict-free or explicitly mediated, and every non-optional evidence
obligation has an authority.  Runtime effectiveness is:
\begin{align*}
effective(A) \iff{}& \forall w \in W_A^{required},\\
                   & evidence(w,A,\ell)\text{ is valid and fresh}.
\end{align*}
Quorum or optional-process semantics must be declared by capability type; the
manager may not silently weaken the quantifier.

\subsection{Threat Model}
We address accidental incompatibility, stale or replayed evidence, conflicting
ownership, provider ambiguity, parent/worker version mismatch, partial rollout,
observer loss, process crashes, and interrupted rollback.  Plugins and
providers may be buggy and may report inconsistent metadata.  We assume the
host, package database, and process identity channel are administratively
trusted.  We do \emph{not} sandbox malicious Python, defend a compromised host,
or claim that signed metadata makes arbitrary plugin code safe.  This boundary
follows the upstream warning that plugins execute with application authority
\cite{vllm-security}.

\subsection{Invariants and Failure Semantics}
\system enforces four invariants.  \textbf{I1, evidence monotonicity:} a stronger
state requires new evidence and cannot be inferred from a weaker state.
\textbf{I2, single explicit ownership:} every exclusive resource has one owner
or a declared mediator.  \textbf{I3, authority locality:} evidence is accepted
only from the component that can observe the claimed event.  \textbf{I4,
rollback preservation:} no new plan destroys the last committed valid plan
before the replacement satisfies its commit gate.

Failures are structured outcomes, not booleans.  A missing provider or
unsupported version makes a plan inadmissible.  A service outage can preserve
\emph{enabled} intent while marking health degraded and refusing a new launch.
A missing worker attestation leaves the activation partial, never globally
effective.  Stale evidence is rejected by launch, process, host, manifest, and
runtime-version identity.  Rollback either restores the recorded predecessor
or reports a bounded operator action; it never invents success.

\section{Design}
\subsection{Typed Static Manifests}
Discovery reads distribution metadata and a static manifest without importing
implementation modules.  The schema distinguishes bundle identity, component
kind, semantic capability, host/API range, provider, required service, resource
claims, and activation/evidence contracts.  Static discovery avoids executing
plugin code merely to decide whether it is eligible.  Namespace ownership is
typed: the project owns \code{vllm\_hust.*}; a provider may delegate to an
official \code{vllm.*} hook but may not mint an upstream-looking group.

The namespace incident becomes a schema fixture: package metadata must contain
exactly the declared group and bundle identity, and a clean installed-package
probe must observe it.  This closes the test gap that allowed a constant to
change while mock-based tests continued to pass.

\subsection{Capability Negotiation}
Capabilities name semantics rather than package names.  A plugin can require a
versioned preemption-policy interface, a connector contract, or process-owned
invocation evidence.  The host and providers advertise concrete capabilities
with provenance.  Negotiation produces a graph whose edges explain which
authority satisfies each requirement.  Missing, ambiguous, or incompatible
edges reject activation.  Untested version tuples remain \emph{unverified}, not
implicitly compatible.

\subsection{Conflict Planning}
Plugins claim resources such as a command-line/configuration key, hook slot,
scheduler policy, connector, service endpoint, or deployment field.  The
planner constructs an ownership graph and detects duplicate exclusive claims,
incompatible values, and conditional conflicts.  For example, two plugins that
both write \code{--kv-transfer-config} conflict even if their package names and
providers differ.  Two plugins sharing a read-only health provider need not
conflict.  The output is a deterministic, reviewable plan rather than an
order-dependent sequence of mutations.

\subsection{Explicit Lifecycle}
The state vector is not a scalar:
\begin{center}
\code{installed, configured, enabled, runtime\_effective}.
\end{center}
Installed is distribution evidence; configured is validated saved input;
enabled is durable operator intent; runtime-effective is launch-scoped evidence
that the selected implementation ran.  Reachability, health, degradation, and
compatibility are orthogonal projections.  Thus an external outage does not
erase intent, and saved intent does not certify a running process.

\subsection{Process-Owned Runtime Evidence}
Each activation plan emits obligations keyed by launch identity, process role,
plugin/component digest, host/runtime version, and capability.  A worker (or an
observer in its authority domain) records load and invocation events.  The
aggregator verifies identity, freshness, expected process membership, and the
declared completion predicate.  Parent-process import is insufficient for a
worker-side hook.  Evidence is compact lifecycle telemetry, not prompt content.

In the running example, the static probe rejects the wrong namespace before
launch.  If static discovery succeeds but one worker never invokes the hook,
the plan may remain enabled and partially loaded, but cannot become globally
runtime-effective.  This same distinction covers inert imports and mixed
worker versions without special-case status strings.

\subsection{Fail-Closed Activation and Deterministic Rollback}
Planning is side-effect free.  Activation proceeds through prepare, launch,
observe, and commit.  Missing required services, contradictory capabilities,
or incomplete evidence stop before commit.  The manager stores the canonical
predecessor plan, rendered inputs, version identities, and authority owners.
Rollback re-renders that predecessor and requires its own evidence gate.
External-service lifecycle remains operator-owned; rollback can restore client
configuration without pretending to downgrade a service it does not control.

\subsection{Authority Separation}
Core owns manifests, intent, planning, and evidence validation.  Providers own
translation and provider-specific checks.  The runtime owns process execution
and invocation facts.  External operators own services, clusters, drivers, and
data.  A provider protocol exposing plan/render/check therefore cannot silently
apply, delete, or restart infrastructure.  This separation limits both failure
blast radius and the claims that status may make.

\section{Implementation}
The prototype is implemented in Python as \code{vllm-hust-ext}.  Core modules
parse manifests, discover installed bundles, maintain user configuration,
resolve compatibility, reject duplicate registrations and resource conflicts,
and dispatch to host providers.  The CLI exposes list, inspect, configure,
enable, check, and run-oriented paths.  Providers currently cover vLLM-native
policies, Mooncake connectors/services, and rendered production-stack inputs.

The implementation already keeps static bundle discovery separate from
implementation entry points and distinguishes persistent intent from provider
health.  It rejects mutating provider plans and documents that only a
process-owned observer can establish runtime effectiveness.  Existing evidence
includes provider-specific integration campaigns, but those campaigns were not
designed as the matched experiment required by this paper.

Two important pieces are incomplete.  First, capability and conflict data need
one frozen machine-readable taxonomy with an independently labeled oracle.
Second, process-owned evidence needs a common launch identity and aggregation
protocol across parent/worker topologies.  Consequently we claim an
implementation foundation and a verified namespace failure, not a completed
end-to-end \system evaluation.

\section{Evaluation Methodology}
The evaluation is organized around four falsifiable hypotheses aligned with
issues M0--M3.  \textbf{H1:} in preregistered failure cells, \system has zero
false-effective outcomes and at least 99\% required-event coverage.
\textbf{H2:} conflict precision and recall each reach 0.95 against an
independent oracle, without rejecting all compositions.  \textbf{H3:} every
unsupported tested tuple fails before serving and at least 95\% of injected
upgrade/rollback failures meet the declared recovery objective or bounded
operator action.  \textbf{H4:} relative to manual integration, added median
startup cost is at most 5\% or 250\,ms, throughput regression at most 1\%, and
p99 request-latency regression at most 2\% in every supported formal cell.

\subsection{Baselines and Workloads}
We compare (1) vanilla vLLM entry-point discovery/loading
\cite{vllm-plugin,vllm-arch}, (2) manual integration using documented flags,
configuration, and operator checks, and (3) \system.  Arms pin identical host,
plugin, provider, model, workload, hardware, process topology, and warm/cold
state.  Formal cells use at least three independent service starts and
alternating arm order.  Agent, repeated-prefix, and online chat workloads cover
different lifecycle and connector behavior; CPU-only fixtures isolate planner
correctness.

\subsection{Matrices and Metrics}
The failure matrix injects missing distributions, malformed manifests,
unsupported versions, provider ambiguity, parent/worker mismatch, inert hooks,
service outage, replayed evidence, partial rollout, process crash, and
interrupted rollback.  The composability matrix covers single plugins,
compatible and incompatible pairs, and three-plugin sets across flags, hooks,
schedulers, connectors, services, and provider translations.  The version
matrix crosses current, previous, and next host versions with below/in/above
plugin ranges, old manifests, and mixed workers.

Primary metrics are activation-event coverage, false-effective rate, conflict
precision/recall, rollback success and recovery time, unsupported-tuple
admission, and cross-process agreement.  Costs include phase-separated startup
time, throughput, latency quantiles, CPU time, memory, process messages, and
evidence bytes.  Raw observations follow the checked-in JSON schema; failures
remain rows rather than being removed from aggregates.

\begin{table}[t]
\centering
\small
\caption{Current evidence and formal-evaluation status.  No planned value is a
result.}
\begin{tabularx}{\columnwidth}{@{}lXl@{}}
\toprule
Claim area & Current evidence & Status \\
\midrule
Static contracts & 66 CPU tests on main & observed \\
Namespace safety & installed entry: expected 1; changed lookup: 0 & negative \\
Conflict P/R & preregistered labeled matrix & \tbd{} \\
False-effective & process failure matrix & \tbd{} \\
Rollback & upgrade/rollback injection & \tbd{} \\
Overhead & matched three-arm experiment & \tbd{} \\
\bottomrule
\end{tabularx}
\end{table}

\subsection{Preliminary Evidence and Missing Results}
On repository main, 66 CPU tests pass.  On commit
\code{3201e133}, those tests also pass despite the namespace change.  A clean
installed-package probe then observes one Mooncake entry point in the documented
group, zero in the changed group, and zero discovered bundles.  This is a real
counterexample to ``unit tests plus installation imply activation,'' and it has
been retained as the M0 negative fixture rather than merged.

All quantitative claims beyond this counterexample are planned.  Table and
figure generators must refuse to replace missing values with zero.  The final
paper will include per-cell confidence intervals and raw failure counts.  Until
H1--H4 are measured, the strongest defensible claim is that the architecture
makes evidence obligations explicit and catches a class of namespace errors
that the previous unit suite missed.

\section{Related Work}
vLLM documents a plugin system and a multi-process engine architecture
\cite{vllm-plugin,vllm-arch}.  \system does not replace those mechanisms; it
adds planning and evidence semantics around them.  Python packaging entry
points support extensible discovery \cite{python-entrypoints}, but package
metadata does not establish runtime invocation.  Dependency solvers reason
about version constraints, while \system also reasons about semantic
capabilities, exclusive runtime resources, and process evidence.

Service meshes, operators, and deployment controllers provide health and
rollout mechanisms.  Our provider boundary deliberately consumes such evidence
without becoming their control plane.  Feature-flag systems separate intent
from rollout but usually do not prove that a specific inference worker invoked
a selected implementation.  Software attestation and provenance systems bind
artifacts to identities; \system uses a narrower, operational notion of
launch-scoped evidence and does not claim cryptographic remote attestation.

Plugin isolation and supply-chain security are complementary.  Because vLLM
plugins execute in application processes, typed manifests are not a sandbox.
Our security contribution is limited to explicit activation authority,
provenance, conflict visibility, and fail-closed state projection.

\section{Discussion and Limitations}
Evidence has cost and can itself fail.  A synchronous per-request proof would
be unacceptable; the design therefore targets lifecycle and sampled invocation
events, but H4 must show that this does not move meaningful work onto the
request path.  Process membership is also subtle under elastic workers.  The
plan must declare whether membership is fixed, epoch-based, or quorum-based;
otherwise a disappearing worker could make success vacuous.

Conflict taxonomies can be incomplete.  Precision and recall are measured
against an independently labeled oracle, but that oracle may miss semantic
interactions.  We mitigate this with explicit unknown/unverified outcomes and
negative cells, not a universal composability claim.  Deterministic rollback
can restore manager-owned inputs, not external data or services; bounded
operator action is therefore a first-class outcome.

The current implementation predates the complete evidence protocol.  Its
provider demonstrations establish feasibility in selected cells, not general
correctness.  The paper should be stopped or reframed as an engineering report
if runtime intent cannot be separated from process effect, conflict thresholds
remain unmet after one correction cycle, rollback loses the predecessor plan,
or overhead exceeds the frozen bounds.

\section{Conclusion}
Inference plugin systems need a truth model stronger than successful
installation or saved enablement.  \system treats activation as a typed,
conflict-checked plan with explicit authorities and evidence obligations.  It
separates lifecycle predicates, requires runtime facts from the processes that
own them, fails closed, and preserves a deterministic rollback target.  A real
namespace regression already demonstrates the value of installed-package
evidence over mock-only tests.  The broader systems claim remains contingent on
the preregistered multi-process, composability, version-skew, recovery, and
overhead evaluation.  This evidence boundary is intentional: the architecture
is useful only if its own status claims are as disciplined as those it imposes
on plugins.

\bibliographystyle{plain}
\bibliography{references}
\end{document}
